\scanpage{47}
\mspage{6}
% L1–2
exactement les $\{\unsure{\zeta|\xi}_i,\zeta'_j\}$ $(i\ne j)$ $\{\zeta_i,\zeta_{ij}\}$ $(i\ne j)$ $\{\zeta'_i,\zeta_{ij}\}$ $(i\ne\unsure{j})$ et $\{\zeta_{ij},\zeta_{k\ell}\}$ \inserted{avec $\{i,j\}\cap\{k,\ell\}=\varnothing$}. \deleted{On note que}

% L3–8
\uline{Lemme} Si $\delta,\delta'$ sont deux sommets (distincts) liés alors $\exists!\,\delta''$ lié à $\delta$ et $\delta'$.

En effet, pour $\{\zeta_i,\zeta'_j\}$ c'est $\zeta_{ij}$, pour $\{\zeta_i,\zeta_{ij}\}$ c'est $\zeta'_j$, pour $\{\zeta'_i,\zeta_{ij}\}$ c'est $\zeta_j$ ; pour $\{\zeta_{ij},\zeta_{k\ell}\}$ c'est $\zeta_{rs}$ où $r,s$ sont les deux entiers $1\le\ \le6$ distincts de $i,j,k,\ell$.
% slanted; clipped at right edge
\margin{\unsure{NB} \illegible{1}}

% L9–14
\uline{Définition} Graphe cubique = graphe isomorphe au graphe \unsure{standart|standard} précédent (construit en prenant \deleted{\unsure{$I$}} $[1,6]\amalg[1,6]\amalg\unsure{\mathfrak P_2|\mathscr P_2}([1,6])$ comme ens de sommets).

À un tel graphe \inserted{$\Delta$ (\unsure{comme : tt graphe|comme à tout graphe|comme : tel graphe})} on associe (pour un anneau comm. $k$ donné) la forme bilinéaire \inserted{sym.} sur $k^{[\Delta]}\times k^{(\Delta)}$ donnée par
\[
\delta\delta'=\begin{cases}
-1&\text{si }\delta=\delta'\\
1&\text{si }\delta\ne\delta',\ \delta\text{ et }\delta'\text{ liés}\\
0&\text{si }\delta\unsure{\ne|=}\delta',\ \delta\text{ et }\delta'\text{ non liés}
\end{cases}
\]
% L15–22
Considérons le noyau \inserted{$N(\Delta)$} de cette forme, et soit $\widehat E_k(\Delta)=k_{\unsure{\zeta|?}}^{[\Delta]}/N(\Delta)$, qui est donc muni d'une forme bilinéaire sym. (non dégénérée au sens faible)

Dans le cas d'un graphe cubique, celle-ci est donc non dégénérée (et $\widehat E_k(\Delta)$ est un module libre de rang 7) comme on voit sur le cas standard.
% archive sheet number: 46
